% ---
% id: reductions/def:joint-family
% kind: def
% title: "Joint alphabet, classes and the joint family of a private setting"
% summary: >-
%   A private setting with exact class weights joins the outcome in one joint
%   alphabet, partitioned into one class per setting, and its conditional
%   POVMs form one joint family with each weight inside the elements. The
%   family satisfies the class-marginal identity and opens no target block.
% prominence: standard
% status: accepted
% uses: [lem:linear, def:exactadm]
% mentions: [ass:probe, ass:tsm, def:reduction-instance, thm:instrument-instance, rem:unsound, rem:intertwining, cor:blockdiag]
% assumes: []
% refs: []
% provenance: newly-derived
% lean:
% notes: "Assumes nothing, like the receiver's def:setting-family: the physical reading of q_t as the measuring party's own weight is A-PROBE and A-TSM, declared by the instance that builds a joint family from a probe. rem:unsound's eq:probeidentity is the identity eq:classmarginal of this definition."
% ---

A receiver with a private per-round setting, such as a probe setting or an
internal coin of known bias, is described by one conditional POVM per value
of the setting, and the analysis reads the pair (outcome, setting) as one
joint outcome with the setting's weight inside the element
(\repodoc{spec/notation.md#26-comma-vs-bar-on-povm-subscripts}). The
definition names that object once and in two layers. The combinatorial layer
is the joint alphabet with its classes and their weights; the abstract
\tref{reduction instance} of Definition~\ref{def:reduction-instance} is stated over
it, and the branching instance of Theorem~\ref{thm:instrument-instance}
extends it by a branch label. The POVM layer is what a device with a private
setting hands this module. The weights are the measuring party's own
(Assumption~\ref{ass:probe}) and enter only through the statistics of the
joint outcome (Assumption~\ref{ass:tsm}). The identity
\eqref{eq:classmarginal} is a fact about the \emph{physical} model; which
squashed form of it is sound is the subject of Remark~\ref{rem:unsound}, and
why the joint alphabet is built \emph{before} the first squashing link, and
what may still be relabelled after one, is Remark~\ref{rem:intertwining}.

\begin{definition}\label{def:joint-family}
Let $\Yset$ be a finite non-empty ordered
\term[The finite ordered set of device outcomes that a private setting is joined to.]{outcome alphabet},
let $\Tset$ be a finite non-empty ordered set of
\term[The finite ordered set of values of a private per-round setting of the receiver.]{settings},
and let $(q_t)_{t\in\Tset}$ be
\term[The declared probabilities of the settings: positive lattice values that sum to one exactly.]{class weights},
each a declared dyadic rational on the declared lattice
(Definition~\ref{def:exactadm}\,(b)), with $q_t>0$ and
$\sum_{t\in\Tset}q_t=1$ \emph{exactly}.
\begin{enumerate}[label=(\alph*),leftmargin=2.4em]
\item \emph{(the joint alphabet, its classes and weights)} The
  \term[The product of the outcome alphabet and the settings, ordered with the outcome varying fastest.]{joint alphabet}
  is $\Vset\coloneqq\Yset\times\Tset$, ordered with $y$ varying
  fastest, which is its \tref{declared order}; its
  \term[The blocks of the joint alphabet, one per setting: every outcome paired with that setting, in the order of the outcome alphabet.]{classes}
  are $\Vset_t\coloneqq\Yset\times\{t\}$, $t\in\Tset$, a partition
  $\Vset=\bigsqcup_{t\in\Tset}\Vset_t$, each class ordered as $\Yset$ is; and
  the triple $\big(\Vset,\{\Vset_t\}_{t\in\Tset},\{q_t\}_{t\in\Tset}\big)$ is
  the \emph{joint alphabet with classes and weights}. When no private setting
  is declared, $\Tset$ is a singleton $\{\ast\}$ with $q_{\ast}=1$, and
  $\Vset=\Yset\times\{\ast\}$ is identified with $\Yset$: one class, of weight
  one.
\item \emph{(the joint family)} Let $\Bsys$ be a system and, for each
  $t\in\Tset$, let $\{\widetilde M_{y|t}\}_{y\in\Yset}$ be a
  \emph{conditional} POVM on $\Bsys$ --- complete at each fixed $t$,
  $\sum_{y}\widetilde M_{y|t}=\idon{\Bsys}$. Their
  \term[The POVM over the joint alphabet whose element at an outcome and a setting is the class weight times the conditional element of that outcome given that setting.]{joint family}
  is
  \begin{equation}\label{eq:jointfamily}
    M_{(y,t)}\;\coloneqq\;q_t\,\widetilde M_{y|t},
    \qquad (y,t)\in\Vset ,
  \end{equation}
  formed from the conditional elements and the weights by one product each,
  exactly. It is a POVM on $\Bsys$ over $\Vset$ by
  Lemma~\ref{lem:linear}\,(ii) --- $\sum_{(y,t)}q_t\widetilde M_{y|t}
  =\sum_tq_t\idon{\Bsys}=\idon{\Bsys}$ --- and it satisfies the
  \term[The operator identity that the elements of each class of a joint family sum to the class weight times the identity.]{class-marginal identity}
  \begin{equation}\label{eq:classmarginal}
    \sum_{v\in\Vset_t}M_v\;=\;q_t\,\idon{\Bsys}
    \qquad\text{for every }t\in\Tset ,
  \end{equation}
  an operator identity on $\Bsys$, by completeness of the conditional POVM at
  $t$.
\end{enumerate}
Two consequences are immediate and are used throughout. First, given
conditional relations $\widetilde M_{|t}\dsq_{\Lambda}\widetilde F_{|t}$,
$t\in\Tset$, with one common channel $\Lambda$, Lemma~\ref{lem:linear}\,(ii)
gives the relation $\{q_t\widetilde M_{y|t}\}\dsq_{\Lambda}\{q_t\widetilde
F_{y|t}\}$ for the joint families: the setting never enters the target space,
no block is opened for it, and the weights are carried by the target elements
themselves. A relation for a joint family need not have a target of this
form. Second, where every $\widetilde M_{y|t}$ is block
diagonal for an orthogonal decomposition of $\Hs{\Bsys}$, so is every
$M_{(y,t)}$, and \eqref{eq:classmarginal} holds block by block; where such a
joint family is itself a target, the decomposition being finite and the
carried operators block diagonal for it as well,
Corollary~\ref{cor:blockdiag} applies to it.
\end{definition}
